\begin{theorem}[Equality of virtual images under a $G$-injective map]%
    \label{thm:peps_g_injective_image_equality}
    \lean{TNLean.PEPS.IsGInjective.averageMap_eq_iff_apply_eq}
    \leanok
    \uses{def:peps_g_injective}
    Let $G$ be finite, let $\rho$ be a complex representation of $G$,
    and let $T$ be $G$-injective. Write
    $\Pi=|G|^{-1}\sum_{g\in G}\rho(g)$ for the invariant projector.
    For virtual vectors $x,y$, one has $Tx=Ty$ if and only if
    $\Pi x=\Pi y$. This is a consequence of
    \cite[Definition~5.1]{Schuch2010PEPS}.
\end{theorem}

\begin{proof}\leanok
    \uses{thm:peps_g_injective_invariants_range}
    The equality $Tx=Ty$ is equivalent to
    $x-y\in\ker T=\ker\Pi$, hence to $\Pi x=\Pi y$.
\end{proof}

\begin{theorem}[A regular $G$-injective site distinguishes translation orbits]%
    \label{thm:peps_regular_site_orbit_detection}
    \lean{TNLean.PEPS.IsGInjective.regularSite_coeff_eq_iff_common_translation}
    \leanok
    \uses{def:peps_g_injective, def:peps_regular_site_gram}
    Let $I$ be a nonempty finite set of incident legs. Let
    $a(\eta)\in\mathcal H$ be the physical vector associated with a
    configuration $\eta\in G^I$, and suppose that the corresponding
    site map is $G$-injective for simultaneous left translation.
    Then $a(\eta)=a(\theta)$ if and only if
    $\eta=g\theta$ for some $g\in G$. This is an auxiliary
    reconstruction consequence of
    \cite[Definition~5.1]{Schuch2010PEPS}.
\end{theorem}

\begin{proof}\leanok
    \uses{thm:peps_g_injective_image_equality,
        def:peps_regular_site_gram, thm:peps_regular_site_gram}
    Write $\delta_\eta$ for the coordinate vector at $\eta$.
    Equality of the physical vectors gives
    $\Pi\delta_\eta=\Pi\delta_\theta$. The matrix entry
    $\Pi_{\eta,\eta}=|G|^{-1}$ is nonzero, whereas
    $\Pi_{\eta,\theta}=0$ if $\eta$ and $\theta$ lie in distinct
    translation orbits. Thus the configurations lie in the same
    orbit. The converse is the virtual invariance of the site map.
\end{proof}

\begin{theorem}[Recovering the group of simultaneous virtual relabellings]%
    \label{thm:peps_regular_site_permutation_recovery}
    \lean{TNLean.PEPS.IsGInjective.regularSite_relabelling_invariant_iff}
    \leanok
    \uses{def:peps_g_injective, def:peps_regular_site_gram}
    In the setting of Theorem~\ref{thm:peps_regular_site_orbit_detection},
    suppose that $I$ has at least two elements. A map $e:G\to G$
    satisfies $a(e\circ\eta)=a(\eta)$ for every configuration $\eta$
    if and only if there exists $g\in G$ with $e(h)=gh$ for every
    $h\in G$. In particular, the simultaneous permutations of
    group labels that preserve the site tensor are exactly the
    left translations.
\end{theorem}

\begin{proof}\leanok
    \uses{thm:peps_regular_site_orbit_detection, thm:peps_regular_site_gram}
    Choose distinct $i,j\in I$ and fix $h\in G$. Let $\eta_j=h$
    and $\eta_k=1$ for $k\ne j$. By
    Theorem~\ref{thm:peps_regular_site_orbit_detection}, there exists
    $g_h\in G$ with $e\circ\eta=g_h\eta$. Evaluation at $i$ gives
    $g_h=e(1)$, and evaluation at $j$ gives $e(h)=e(1)h$.
    Conversely, each left translation preserves the site tensor by
    its virtual invariance.
\end{proof}
